% !TEX TS-program = pdflatex
% !TEX encoding = UTF-8 Unicode

% Example of the Memoir class, an alternative to the default LaTeX classes such as article and book, with many added features built into the class itself.

%\documentclass[12pt,a4paper]{memoir} % for a long document
\documentclass[12pt,a4paper,article]{memoir} % for a short document

\usepackage[utf8]{inputenc} % set input encoding to utf8
%LANGUAGE FR
\usepackage[cyr]{aeguill}
\usepackage[francais]{babel}
\usepackage{layout}
\usepackage[left=2.5cm, right=2.5cm]{geometry}
\usepackage{lipsum}
\usepackage{graphicx}
\graphicspath{{images/}}

% Don't forget to read the Memoir manual: memman.pdf

%%% Examples of Memoir customization
%%% enable, disable or adjust these as desired

%%% PAGE DIMENSIONS
% Set up the paper to be as close as possible to both A4 & letter:
\settrimmedsize{11in}{210mm}{*} % letter = 11in tall; a4 = 210mm wide
\setlength{\trimtop}{0pt}
\setlength{\trimedge}{\stockwidth}
\addtolength{\trimedge}{-\paperwidth}
%\settypeblocksize{*}{\lxvchars}{1.618} % we want to the text block to have golden proportionals
\setulmargins{40pt}{*}{*} % 50pt upper margins
%\setlrmargins{*}{*}{1.618} % golden ratio again for left/right margins
%\setheaderspaces{*}{*}{1.618}
%\checkandfixthelayout 
% This is from memman.pdf

%%% \maketitle CUSTOMISATION
% For more than trivial changes, you may as well do it yourself in a titlepage environment
\pretitle{\begin{center}\sffamily\huge\MakeUppercase}
\posttitle{\par\end{center}\vskip 0.5em}

%%% ToC (table of contents) APPEARANCE
\maxtocdepth{subsection} % include subsections
\renewcommand{\cftchapterpagefont}{}
\renewcommand{\cftchapterfont}{}     % no bold!

%%% HEADERS & FOOTERS
\pagestyle{ruled} % try also: empty , plain , headings , ruled , Ruled , companion

%%% CHAPTERS
\chapterstyle{hangnum} % try also: default , section , hangnum , companion , article, demo

\renewcommand{\chaptitlefont}{\Huge\sffamily\raggedright} % set sans serif chapter title font
\renewcommand{\chapnumfont}{\Huge\sffamily\raggedright} % set sans serif chapter number font

%%% SECTIONS
\hangsecnum % hang the section numbers into the margin to match \chapterstyle{hangnum}
\maxsecnumdepth{subsection} % number subsections

\setsecheadstyle{\Large\sffamily\raggedright} % set sans serif section font
\setsubsecheadstyle{\large\sffamily\raggedright} % set sans serif subsection font

%% END Memoir customization

\title{Master Degree Dissertation}
\author{Étudiant: Cyril Annette\\
Responsable: Pierre Bret\\
Entreprise: Smile Montpellier\\
Du 02 avril au 27 septembre 2013\\
}
\date{} % Delete this line to display the current date



%%% BEGIN DOCUMENT
\begin{document}


\maketitle
\begin{center}
	\includegraphics[width=300px]{smile.jpg}
	\includegraphics[width=300px]{supinfo.jpg}
\end{center}

\newpage
\tableofcontents* % the asterisk means that the contents itself isn't put into the ToC
\newpage

\chapter{Introduction}
\section{Présentation du candidat}

\paragraph{}
Avant de vous aventurer plus loin, cher lecteur, il convient d'en apprendre plus sur "le candidat". C'est dans cette introduction, par ces quelques lignes, que j'ai décidé de me présenter à vous de façon différée.

\paragraph{}
Je m'appelle Cyril ANNETTE, et ce mémoire sera un de mes nombreux accomplissements, maintenant que j'arrive au terme d'un cursus de cinq ans visant à me muer en un ingénieur digne de ce titre. Dans celui-ci, vous découvrirez par quels moyens et dans quelles circonstances j'ai effectué la transition entre la couche scolaire et celle du  vrai monde: le travail en entreprise.

\paragraph{}
Complexe tâche que celle de se présenter soi-même dans un document écrit. Aussi, j'essaierai d'être le moins ennuyeux et complaisant possible. 
C'est en 2008 que j'ai commencé mon initiation à l'informatique. Je sortais d'un baccalauréat scientifique, alors jeune et innocent, et motivé par rien de plus que cette sensation d'aise que j'avais en manipulant un ordinateur. A cette époque je n'imaginais pas toutes les branches en lesquelles se décomposait les Technologies de l'Information, ni la richesse et la complexité de chacune d'entre elles. 

\paragraph{Mon choix d'orientation \\}
C'est pourtant vers l'une d'entre elles que je suis rapidément dirigé. L'optique de la création m'a toujours intéressé car elle donne corps aux idées les plus originales. C'est donc en tant que développeur que j'abordais mon cursus. 
Mais comme précisé plus haut, l'Informatique avec un "I" diverge toujours en plus de branches qu'on ne l'imagine. Dans la famille du développement, je fis un nouveau choix pour affiner mon angle d'expertise. 
Le monde du Web est en endroit fascinant. Au delà de la vision que Monsieur Tout-le-monde peut en avoir, les technologies du web offrent un environnement en évolution permanente, mais profondément marqué par les aieux qui lui ont ouvert la voix. Ainsi, on peut se voir rapidement devoir choisir entre des technologies qui ont déjà fait leurs preuves (tel que PHP/MySQL) ou des technologies émergentes mais à fort potentiel (tels que les Node.js et MongoDB). Des profils radicalement différents se forment donc au sein de cette spécialité pourtant déjà très spécifique. Si on ajoute à tout celà les solutions clefs-en-main qui proposent un ensemble de technologies, le nombre de combinaisons possible devient impressionnant.
Tout le monde a donc sa carte à jouer, et les technologies du web offrent de nombreuses opportunités d'embauches. 


\section{Presentation de l'entreprise}

\paragraph{} \lipsum

\subsection{test}

\paragraph{} \lipsum

\subsubsection{subtest}

\paragraph{} \lipsum

\newpage
\chapter{Projets}

\section{Analyse du contexte (RENAME ME) }

\paragraph{} \lipsum

\section{Problématiques (RENAME ME) }

\paragraph{} \lipsum

\section{Méthodes habituellement utilisées pour une situation présentant des similitudes (RENAME ME) }

\paragraph{} \lipsum

\newpage
\chapter{Prise de décisions}

\section{Exposé des décisions prises et des interventions menées par le stagiaire pour résoudre le problème (RENAME ME) }

\paragraph{} \lipsum

\section{Démonstration d’une originalité dans l’élaboration et la mise en œuvre de la solution (RENAME ME) }

\paragraph{} \lipsum

\section{ Analyse de l’approche choisie (RENAME ME) }

\paragraph{} \lipsum

\newpage
\chapter{Conclusions}

\section{ Réflexion sur le stage et le mémoire (RENAME ME) }

\paragraph{} \lipsum

\section{ Conclusion }

\paragraph{} \lipsum

\section{ Annexes }

\paragraph{} \lipsum

\end{document}
